\subsection{Regularity on the exceptional torus}\label{sec:cartan-torus-regularity}
The field calculation extends to the local rings on the exceptional torus and gives an integral obstruction to specializing this Cartan construction.

\begin{proposition}\label{prop:cartan-torus-regularity}
Use the notation of Proposition~\ref{prop:divisorial-full-cartan}. In the fraction field of the regular full-level ring $R_n$, put
\[
 \pi=x^{1-p}y^p,\qquad t=x^B/y^A,\qquad
 S_0=R_n[\pi,t,t^{-1}].
\]
For $c\in k^*$, let $[c]\in W(k)$ be its Teichm\"uller representative and set
\[
 B_c=(S_0)_{(\pi,t-[c])},\qquad A_c=B_c^C,
 \qquad P=p^{n-1}.
\]
Then $B_c$ is a regular local ring of dimension two, the principal Cartan acts on $B_c$, and $A_c$ is regular. The unique reduced vertical divisor of $A_c$ is regular and Cartier. It has a local equation $\rho$ for which
\[
 \rho/\pi^P\in B_c^*,\qquad
 A_c/(\rho)=k[t^P,t^{-P}]_{(t^P-c^P)}
\]
as subrings of the divisorial residue field. After multiplication by a unit of $A_c$, one may arrange $\operatorname{in}(\rho)=\operatorname{in}(\pi)^P$. Moreover
\begin{equation}\label{eq:cartan-torus-p-factorization}
 p=u\rho^{(p^4-1)p^{n-1}},\qquad u\in A_c^*.
\end{equation}
\end{proposition}

\begin{proof}
\par\medskip\noindent\textbf{The regular chart and the action.}\enspace
The identity $(1-p)A+pB=1$ gives the inverse monomial formulas
\begin{equation}\label{eq:cartan-torus-coordinates}
 x=\pi^A t^p,\qquad y=\pi^B t^{p-1}.
\end{equation}
The quotient $S_0/(\pi)$ is exactly $k[t,t^{-1}]$. Indeed $x,y$ vanish in this quotient by~\eqref{eq:cartan-torus-coordinates}, and they generate the maximal ideal of $R_n$. Thus the quotient is generated by $t,t^{-1}$ over $k$. The monomial valuation sends $t$ to a transcendental residue in $k(t)$, which excludes any further Laurent-polynomial relation. Consequently $(\pi,t-[c])$ is a maximal ideal, and the prime chain
\[
 (0)\subsetneq(\pi)\subsetneq(\pi,t-[c])
\]
shows that its localization has dimension at least two. Its maximal ideal has two generators, so Krull's height theorem gives the reverse inequality. This proves regularity of $B_c$, with parameters $\pi,t-[c]$, and
\begin{equation}\label{eq:cartan-torus-original-divisor}
 B_c/(\pi)=k[t,t^{-1}]_{(t-c)}.
\end{equation}
These are uncompleted local rings; their fraction field remains $\operatorname{Frac}(R_n)$.

For every $g\in C$, the principal level action and the torsion initial forms proved above give
\[
 w(gx-x)>A,\qquad w(gy-y)>B.
\]
In the regular local ring $B_c$, divisibility by $\pi^a$ is equivalent to divisorial value at least $a$. One can verify this by intersecting the height-one valuation rings: $\pi$ is a unit at every other height-one prime. Since $gx-x$ and $gy-y$ belong to $R_n\subset B_c$, formula~\eqref{eq:cartan-torus-coordinates} implies
\[
 gx/x,\ gy/y\in1+\pi B_c.
\]
The monomial definitions then give $g\pi/\pi,gt/t\in1+\pi B_c$. Thus $g$ sends the generators of $S_0$ into $B_c$, fixes its reduction modulo $\pi$, and preserves the closed point $t=c$. Denominators outside that point map to units. Applying the same argument to $g^{-1}$ proves that $g$ is an automorphism of $B_c$.

All invariant rings below are finite Noetherian local quotients. Here is a direct verification. The full-level ring is finite over the fixed Noetherian deformation base $R_0$. Form the $R_0$-algebra $S$ generated by all translates of a finite set of $R_0$-algebra generators of $S_0$. Then $S$ is a finite-type $C$-stable subalgebra of $B_c$. If $\mathfrak m$ is the contraction of the maximal ideal of $B_c$, the inclusions $S_0\subset S\subset B_c$ give $B_c=S_{\mathfrak m}$. For any subgroup $H$, the coefficients of the finitely many orbit polynomials generate a Noetherian subring of $S^H$ over which $S$ is finite. Its submodule $S^H$ is therefore finite and Noetherian. Since $\mathfrak m$ is $H$-stable, the norm of any denominator outside $\mathfrak m$ is invariant and remains outside it. Hence
\[
 (B_c)^H=(S^H)_{\mathfrak m\cap S^H},
\]
and $B_c$ is finite over this ring. Normality follows from normality of $B_c$. All closed residue fields are $k$, because the action fixes $W(k)$.

\par\medskip\noindent\textbf{The unique vertical divisor.}\enspace
Write $N=p^{2n-2}$ and $E=(p^4-1)N$. The full level structure gives a Weierstrass factorization of $[p](T)$. Differentiating it at the zero point expresses $p$ as a unit times the product of all nonzero first-level torsion coordinates. There are $p-1$ points with initial form a nonzero multiple of $X^N$, and $p(p-1)$ with initial form a nonzero multiple of $Y^N$. Thus
\[
 \operatorname{in}(p)=a_0X^{(p-1)N}Y^{p(p-1)N},
 \qquad a_0\in k^*.
\]
By~\eqref{eq:cartan-torus-coordinates}, this is
\[
 a_0\operatorname{in}(\pi)^E t^{p^2(p-1)N}.
\]
The Laurent coefficient is a unit at every $c\in k^*$. Divisibility in $B_c$ therefore proves
\begin{equation}\label{eq:cartan-torus-upstairs-unit}
 p=\pi^E h,\qquad h\in B_c^*.
\end{equation}
There is exactly one vertical prime through this chart, and the same holds in every intermediate finite invariant quotient.

\par\medskip\noindent\textbf{Field information and actual local rings.}\enspace
We use the complete induction of Proposition~\ref{prop:divisorial-full-cartan}, including its first derivations and actual fixed-field degrees. Its order is cross first, then scalar. The completed fixed fields also describe the divisorial completions of the invariant fraction fields of the local rings. To verify this point, approximate an invariant element of $F$ by an element of $\operatorname{Frac}(R_n)$ and average over the finite subgroup. Division by its order loses only a fixed finite amount of valuation. Thus invariant elements of the dense fraction field are dense in the completed fixed field. This argument takes place in fields and makes no claim about integral averaging.

At the beginning of a complete congruence layer write
\[
 r=p^m,\quad j=n-m-1\ge1,\quad Q=p^{2j},\quad
 U=X^r,\quad V=Y^r,
\]
and introduce the homogeneous quantities
\begin{equation}\label{eq:cartan-torus-layer-coordinates}
 P_0=U^{1-p}V^p=\operatorname{in}(\pi)^r,
 \qquad T=U^B/V^A=t^r.
\end{equation}
They satisfy $U=P_0^A T^p$ and $V=P_0^B T^{p-1}$. The current graded field is $k(U,V)$, its value group is $r\Z$, and its residue field is $k(T)$.

Inductively, let $\mathcal B$ be the current actual local invariant ring. We prove simultaneously that $\mathcal B$ and its reduced vertical divisor are regular. If $f$ generates the vertical prime, its pullback to $B_c$ is $\pi^r$ times a unit: the unique vertical divisor has that multiplicity by the field ramification calculation. Thus
\[
 \operatorname{in}(f)=P_0 a(T),
\]
where $a$ is regular and nonzero at $T=c^r$. The divisor ring is the local ring of $k(T)$ at $c^r$. Indeed it is a discrete valuation ring with that fraction field, dominated by the finite local ring~\eqref{eq:cartan-torus-original-divisor}; the restricted valuation is precisely the valuation at $T=c^r$. A parameter on this curve therefore has residue $g(T)$ with a simple zero at $c^r$, and $g'(c^r)\ne0$. It lifts to an element of $\mathcal B$. These local induction assertions hold at $m=0$ by the explicit chart construction.

\par\medskip\noindent\textbf{The cross step and its integral parameter.}\enspace
By~\eqref{eq:full-cartan-cross-derivation}, the first derivation of $\tau_j$ on $k(U,V)$ is $V^Q\partial_U$, with shift $r(QB-A)$. As $B=p^2$ is zero in characteristic $p$, it satisfies
\begin{equation}\label{eq:cartan-torus-cross-local}
 D_\tau(T)=0,\qquad
 D_\tau(P_0)=(1-p)P_0^{QB-A+1}T^{(p-1)Q-p}.
\end{equation}
Put $d=QB-A>0$. The second expression has a nonzero Laurent-monomial coefficient at the closed point. Since $D_\tau$ kills $a(T)$, the difference $\tau_j(f)-f$ is $f^{d+1}$ times a unit in $\mathcal B$.

For any element of positive divisorial value in $\mathcal B$, the global shift bound gives displacement at least $r(d+1)$. For an element of value zero, the initial form belongs to $k(T)$ and is killed by~\eqref{eq:cartan-torus-cross-local}. Its displacement has value strictly greater than $rd$, hence at least $r(d+1)$ by discreteness. Divisibility along the prime parameter $f$ now shows that the entire augmentation ideal is
\[
 I_{\tau_j}=(\tau_j(b)-b:b\in\mathcal B)=(f^{d+1}).
\]
The principal-augmentation theorem of \citet[Theorem~2 and Proposition~5, pp.~64--69]{KiralyLutkebohmert2013} proves that $\mathcal A=\mathcal B^{\langle\tau_j\rangle}$ is regular, that $\mathcal B$ is finite free over $\mathcal A$, and, since the augmentation ideal is generated by $\tau_j(f)-f$, that
\[
 \mathcal B=\mathcal A[f],\qquad \mathcal B\simeq\mathcal A[Z]/(P_f(Z)),
\]
where $P_f$ is the monic degree-$p$ orbit polynomial and $1,f,\ldots,f^{p-1}$ is an $\mathcal A$-basis. All the hypotheses of that theorem hold here: the rings are Noetherian normal local domains, $\mathcal B$ is regular, and the action has order $p$.

The actual fixed graded field is $k(U^p,V)$ by~\eqref{eq:full-cartan-cross-field}. Its value group is $rp\Z$, and its residue field remains $k(T)$. Therefore $q=\operatorname{Norm}(f)$ has multiplicity one at the unique vertical prime of $\mathcal A$. It has no other prime divisor, since each conjugate of $f$ is an equation of the vertical divisor. Regularity of $\mathcal A$ implies that $q$ generates that prime up to a unit. The monogenic presentation gives the exact isomorphism
\begin{equation}\label{eq:cartan-torus-norm-divisor}
 \mathcal B/(f)\simeq\mathcal A/(q).
\end{equation}
The reduced divisor of $\mathcal A$ is consequently regular. Moreover
\[
 \operatorname{in}(q)=P_0^p a(T)^p,
\]
because all conjugates of $f$ have the same initial form. This constructs the needed parameter inside the actual local invariant ring, without extending a fractional controlled field lift across the closed point.

\par\medskip\noindent\textbf{The scalar step and descent of the reduced divisor.}\enspace
Equations~\eqref{eq:full-cartan-scalar-b}--\eqref{eq:full-cartan-scalar-a} in the field induction show that on $k(U^p,V)$ the first scalar derivation is $D_\sigma=V^Q\partial_V$, with shift $rB(Q-1)$. It kills $P_0^p$ and satisfies
\begin{equation}\label{eq:cartan-torus-scalar-local}
 D_\sigma(T)=-A P_0^{B(Q-1)}T^{1+(p-1)(Q-1)}.
\end{equation}
Relative to the current value group $rp\Z$, its shift is the positive integer
\[
 d_s=B(Q-1)/p=p(Q-1).
\]
Choose a lift $z\in\mathcal A$ of a parameter $g(T)$ of the regular divisor in~\eqref{eq:cartan-torus-norm-divisor}. Its simple zero gives $g'(c^r)\ne0$. Formula~\eqref{eq:cartan-torus-scalar-local} shows that $\sigma_j(z)-z$ is $q^{d_s}$ times a unit: its normalized leading coefficient is
\[
 \frac{-A g'(T)T^{1+(p-1)(Q-1)}}{a(T)^{p d_s}},
\]
which is nonzero at $c^r$. The global shift bound puts every local difference in $(q^{d_s})$. Hence the augmentation ideal is exactly this principal ideal.

The same principal-augmentation theorem proves that $\mathcal A'=\mathcal A^{\langle\sigma_j\rangle}$ is regular and $\mathcal A$ is finite free over $\mathcal A'$. The actual scalar field step has ramification index one and residue field $k(T^p)$, by~\eqref{eq:full-cartan-next-field}. Let $q'$ generate the unique vertical prime of the regular ring $\mathcal A'$. Its pullback is $q$ times a unit, because that prime has ramification index one and no other vertical prime lies above it. Thus $\mathcal A/(q')=\mathcal A/(q)$ is regular, and
\[
 \mathcal A'/(q')\longrightarrow\mathcal A/(q')
\]
is faithfully flat. Regularity descends under this map, so the reduced divisor of $\mathcal A'$ is regular as well as Cartier. Its local ring is
\[
 k[T^p,T^{-p}]_{(T^p-c^{rp})}.
\]
Its vertical parameter pulls back to $\pi^{rp}$ times a unit. In the graded field $k(U^p,V^p)$ its initial coefficient is a rational function of $T^p$ that is regular and nonzero at the indicated point. We have obtained every local induction assertion with $r$ replaced by $rp$.

\par\medskip\noindent\textbf{The final unit factorization.}\enspace
Iteration through all $n-1$ congruence layers proves regularity of $A_c$ and of its reduced vertical divisor. With $P=p^{n-1}$, a local equation $\rho$ has pullback $\pi^P$ times a unit and the asserted divisor ring. Its initial coefficient belongs to $k(t^P)$ and is a unit in that divisor ring. Lifting its inverse to a unit of $A_c$ normalizes $\operatorname{in}(\rho)=\operatorname{in}(\pi)^P$.

Equation~\eqref{eq:cartan-torus-upstairs-unit} and the pullback-unit statement show that $p/\rho^{E/P}$ and its inverse are units in $B_c$. They are invariant, so they belong to $A_c^*$. Since $E/P=(p^4-1)p^{n-1}$, this proves~\eqref{eq:cartan-torus-p-factorization}.
\end{proof}

\begin{corollary}\label{cor:cartan-torus-specialization}
Every local homomorphism $A_c\to\mathcal O_L$ to a mixed-characteristic discrete valuation ring, with $p$ nonzero, satisfies
\[
 (p^4-1)p^{n-1}\mid e_L.
\]
For $p=5$ and $n=4$, this gives $78000\mid e_L$. If in addition $L/K$ is a finite extension with $e(L/K)\le12$, then $e_K\ge6500$.
\end{corollary}
\begin{proof}
Under the local map, $\rho$ has positive integral valuation and the unit in~\eqref{eq:cartan-torus-p-factorization} remains a unit. Therefore
\[
 e_L=(p^4-1)p^{n-1}v_L(\rho).
\]
At $p=5,n=4$, the coefficient is $624\cdot125=78000$. Multiplicativity $e_L=e(L/K)e_K$ gives the last inequality. Existence of such an extension is not asserted.
\end{proof}

The proposition concerns the points $t\ne0,\infty$ of this particular Cartan chart. It supplies no reduction of every surviving torsion comparison to these local rings. The general first-boundary occurrence question is therefore unchanged. The known non-Cartier quotient below explains why a field ramification calculation, without such a local regularity proof, does not imply the same specialization divisibility for arbitrary normal models.
